\documentclass[11pt]{article}

\usepackage{fontspec}
\usepackage{xeCJK}
\usepackage{xcolor}
\usepackage{graphicx}
\usepackage{array}
\usepackage{tabularx}
\usepackage{enumitem}
\usepackage{titlesec}
\usepackage{needspace}
\usepackage{fontawesome}
\usepackage{hyperref}
\CJKsetecglue{\hskip 0.08em plus 0.04em minus 0.02em}

\definecolor{CVBlue}{RGB}{32,58,86}
\definecolor{CVMuted}{RGB}{85,91,99}
\hypersetup{colorlinks=true,urlcolor=CVBlue,linkcolor=CVBlue}

% Use the bundled fonts; slanted metadata uses the regular face.
\setmainfont[
    Path = fonts/Main/,
    UprightFont = texgyretermes-regular.otf,
    BoldFont = texgyretermes-bold.otf,
    ItalicFont = texgyretermes-regular.otf,
    ItalicFeatures = {FakeSlant=0.18},
]{texgyretermes}
% Prefer a Song-style body face; keep a portable bundled fallback.
\IfFontExistsTF{FandolSong-Regular.otf}{
    \setCJKmainfont[BoldFont=FandolSong-Bold.otf]{FandolSong-Regular.otf}
}{
    \IfFontExistsTF{Songti SC}{
        \setCJKmainfont[BoldFont=Songti SC Bold]{Songti SC}
    }{
        \setCJKmainfont[
            Path = fonts/hansans/,
            UprightFont = NotoSansSC-Regular.ttf,
            BoldFont = NotoSansSC-Bold.ttf,
        ]{NotoSansSC}
    }
}
\setCJKfamilyfont{cvheading}[
    Path = fonts/hansans/,
    UprightFont = NotoSansSC-Regular.ttf,
    BoldFont = NotoSansSC-Bold.ttf,
]{NotoSansSC}

\usepackage[
    a4paper,
    left=1.5cm,
    right=1.5cm,
    top=1.3cm,
    bottom=1.3cm,
    nohead
]{geometry}

\setlength{\parindent}{0pt}
\setlength{\parskip}{2pt}
\renewcommand{\baselinestretch}{1.12}
\setlength{\emergencystretch}{1.5em}
\clubpenalty=10000
\widowpenalty=10000
\raggedbottom

\setlist[itemize]{
    leftmargin=1.3em,
    labelsep=0.45em,
    itemsep=1.5pt,
    parsep=0pt,
    topsep=3pt,
    partopsep=0pt
}

\titleformat{\section}
{\large\CJKfamily{cvheading}\bfseries\color{CVBlue}\raggedright}
{}{0em}{}
[{\color{CVBlue}\titlerule[0.4pt]}]
\titlespacing*{\section}{0pt}{11pt}{5pt}

% A fixed date column keeps long headings aligned and allows them to wrap.
\newcommand{\cvheading}[2]{%
    \par\noindent
    \begin{tabularx}{\linewidth}{@{}>{\raggedright\arraybackslash}X@{\hspace{0.8em}}>{\raggedleft\arraybackslash}p{3.4cm}@{}}
        {\CJKfamily{cvheading}\bfseries #1} & {\small\color{CVMuted}#2}
    \end{tabularx}\par\nobreak\vspace{1pt}%
}
\newcommand{\projectheading}[2]{%
    \Needspace{5\baselineskip}%
    \cvheading{#1}{#2}%
}
\newcommand{\affiliation}[1]{%
    {\small\color{CVMuted}#1\par}\nobreak
}
\newcommand{\projectsep}{\par\addvspace{8pt}}
\newcommand{\publication}[3]{%
    \par\Needspace{5\baselineskip}%
    {\bfseries #1\par}\nobreak
    {\small #2\par}\nobreak
    {\small\color{CVMuted}#3\par}%
}

\begin{document}
\pagenumbering{gobble}

% Crop only the blank margins in the logo; keep the image file unchanged.
\noindent
\begin{minipage}[c]{0.22\linewidth}
    \includegraphics[width=3.5cm,trim=90bp 685bp 95bp 700bp,clip]{zju.png}
\end{minipage}\hfill
\begin{minipage}[c]{0.75\linewidth}
    \raggedleft
    {\fontsize{23}{26}\selectfont\CJKfamily{cvheading}\bfseries 王洁睿\par}
    \vspace{3pt}
    {\fontsize{9.5}{12}\selectfont\color{CVMuted}
    \faPhone\ +86 18516265729 \quad +1 412 912 5930\par
    \faEnvelopeO\ \href{mailto:jieruiwang@zju.edu.cn}{jieruiwang@zju.edu.cn} \quad \href{mailto:jieruiw@andrew.cmu.edu}{jieruiw@andrew.cmu.edu}\par}
\end{minipage}
\par

\section{教育背景}
\cvheading{浙江大学，竺可桢学院}{2023.09 -- 至今}
\noindent
\begin{tabularx}{\linewidth}{@{}>{\raggedright\arraybackslash}X@{\hspace{0.8em}}r@{}}
    计算机科学与技术专业，图灵班，本科 & {\small\color{CVMuted}预计毕业时间：2027年6月}
\end{tabularx}\par
\begin{itemize}
    \item GPA：4.2/4.3 \qquad 排名：1/99 \qquad 雅思：7.5
    \item 部分课程：计算机组成与设计（100）、机器学习（98）、操作系统（98）、编译原理（97）、数学分析（H）（96）、线性代数（H）（96）
\end{itemize}

\section{论文}
\publication
    {AuraForge: Beyond Human-Written Tests for Training Secure Vibe Coding Agents}
    {Danqing Wang, Songwen Zhao, Harsh Sharma, \textbf{Jierui Wang}, Andre Vicente Duarte, Ivan Bercovich, Yangruibo Ding, Lei Li.}
    {\textbf{ICLR 在投}}

\projectsep

\publication
    {SeekJudge: A Practical Reward Framework for Reinforcement Learning in Computer-Use Agents}
    {Yang Wan, Zhenhao Zhang, \textbf{Jierui Wang}, Linchao Zhu.}
    {\textbf{ICLR 在投}\par\nobreak
    \href{https://arxiv.org/abs/2607.23263}{https://arxiv.org/abs/2607.23263}\quad\href{https://github.com/ZJUSCL/SeekJudge}{https://github.com/ZJUSCL/SeekJudge}}

\section{科研实习}
\cvheading{卡内基梅隆大学，Li Lab}{2026.06 -- 2026.09}
科研实习生，指导老师：Lei Li

\projectsep

\cvheading{浙江大学，ZJU Scaling Logger（ZJUSCL）}{2025.05 -- 2026.05}
科研实习生，指导老师：朱霖潮

\section{研究与项目经历}

\projectheading{AuraForge: Beyond Human-Written Tests for Training Secure Vibe Coding Agents}{2026.06 -- 至今}
\affiliation{Li Lab}
\begin{itemize}
    \item \textbf{研究目标：}构建可规模化的可执行监督信号，训练代码智能体满足\textbf{隐式安全需求}，从功能行为与特定漏洞对应的安全属性两方面评估生成补丁。
    \item \textbf{任务构造：}主导\textbf{JavaScript/TypeScript 任务构造流水线}，将真实漏洞修复转换为可执行的代码智能体任务。设计\textbf{语言感知的功能实现遮蔽机制}，移除存在漏洞的实现，同时保留独立重实现所需的声明、导出、类型签名与接口。
    \item \textbf{执行基础设施与规模：}重建\textbf{仓库专用的 Node.js 环境}，涵盖运行时、包管理器、工作区、依赖和构建配置；开发\textbf{统一异构 JS/TS 测试运行器的适配器}，支持跨仓库状态的一致功能与安全验证。将\textbf{863 个经筛选的 JS/TS 候选转换为 229 个任务实例}，用于安全代码智能体的评估与训练。
\end{itemize}

\projectsep

\projectheading{SeekJudge：计算机操作智能体的奖励建模}{2026.03 -- 2026.05}
\affiliation{ZJU Scaling Logger（ZJUSCL）}
\begin{itemize}
    \item \textbf{研究目标：}为长程计算机操作智能体开发基于模型的奖励，使在线强化学习无需依赖任务专用的规则验证器。
    \item \textbf{研究方法：}共同开发\textbf{Condense--Ground--Seek--Analyze}框架，将长截图与操作轨迹压缩为语义时间线，将底层操作与对应界面元素及语义关联，并通过 Seek--Analyze 循环迭代定位、检查决定性视觉证据。最终将多智能体评判过程蒸馏至统一的\textbf{9B 评判模型}。
\end{itemize}

\projectsep

\projectheading{TradeActor：面向投资组合配置的多模态市场推理}{2026.03 -- 2026.05}
\begin{itemize}
    \item 开发\textbf{多模态交易框架}，结合股价图表、技术指标图表与市场上下文，使用\textbf{冻结参数的视觉语言模型}提取语义信号，为下游投资组合决策提供支持。
    \item 设计轻量级\textbf{Qwen-0.6B 投资组合策略模型}，将市场上下文编码至指定的\textbf{预测 token}，并通过\textbf{MLP 动作头}将该 token 的最终隐藏状态映射为连续的投资组合配置。
    \item 以\textbf{考虑交易成本惩罚的投资组合收益}为目标优化策略，并通过\textbf{StockBench 端到端回测}评估每日资产配置决策。
\end{itemize}

\projectsep

\projectheading{TimeThinkingLLM：多模态时间序列推理}{2025.09 -- 2026.01}
\begin{itemize}
    \item 开发\textbf{多模态框架}，将\textbf{连续时间序列信号}与\textbf{Qwen-3 的表示空间}对齐，使预训练语言模型能够直接理解时序数据。
    \item 设计\textbf{时间编码器}与\textbf{投影机制}，将时间序列信号转换为与大语言模型\textbf{语义嵌入空间}兼容的表示。
    \item 使用\textbf{基于 GRPO 的强化学习}训练模型进行\textbf{时序信号推理}，并在\textbf{时间序列基准}上评估其时序推理能力。
\end{itemize}

\end{document}
